\begin{abstract}

Self-supervised EEG representation learning has converged on foundation-scale recipes~\cite{labram2024,eegpt2024,neurolm2025,laya2026} (tens of thousands of hours, 5\textendash 10\,M$+$ parameters, multi-GPU). The compact, single-corpus regime in which most laboratories and on-device applications operate is uncharacterized, and standard cross-subject probes can be silently inflated by subject leakage. We address both gaps using \textbf{EEG-LeJEPA}, a 2.86\,M-parameter Joint-Embedding Predictive Architecture trained with Sketched Isotropic Gaussian Regularization (SIGReg)~\cite{lejepa2025} as a leakage-clean testbed. Pretrained for 10{,}000 steps on 50 EEGMMIDB subjects (51--100, disjoint from evaluation) in under 15 minutes on a single consumer GPU, it reaches 68.6\,\%/0.751\,AUC on rest-versus-activity and 62.1\,\%/0.690\,AUC on left-versus-right motor imagery via a frozen-encoder linear probe (LOSO; subjects 1--20 strictly held out); both gains over a random-initialized identical architecture are large and significant ($+11.7$ and $+10.9$\,pp; Benjamini--Hochberg $q<0.002$). Three transferable findings emerge. First, under the disjoint split, predictor-derived feature sources lose their apparent advantage and the simple encoder mean-pool is the strongest and most trustworthy source. Second, the SIGReg projection count must scale with embedding width: an isolation run shows a 2.4$\times$ larger model collapses 8.3\,pp on left-vs-right MI when only $M$ is reduced; at properly-scaled $M$ it is modestly but consistently better across three seeds ($+1.3$ to $+2.0$\,pp). Third, channel-padded cross-dataset transfer outperforms channel-matched re-pretraining ($+4.9$ vs.\ $+1.5$\,pp). A SIGReg-weight ablation confirms complete representational collapse at $\lambda{=}0$; the SSL$+$probe pipeline matches supervised-from-scratch on left/right MI ($-0.2$\,pp, n.s.) and exceeds it on rest/activity ($+8.4$\,pp, $p<0.001$).
\end{abstract}
